% sec_gram_space.tex -- Mode Gram kernel and RKHS completion
\section{The Mode Gram Kernel}%
\label{app:gram-space}

This appendix computes the Gram kernel attached to the comparison kernel
$R_\varepsilon$ and identifies the resulting RKHS.  The closed-form kernel
$K(a,b)=2/(4+(a-b)^2)$ supports the entropy story directly: it encodes
the $L^2(\omega)$ geometry of the mode fluctuation profiles $\delta_t$.
The RKHS framework follows the general analytic translation-invariant
theory of \cite{SunZhou2008}; orthonormal expansions for the Cauchy/Lorentz
class are in \cite{TronarpKarvonen2024}.

\subsection{The Gram kernel computation}

Let
\[
\omega(\dd y):=\frac{\dd y}{\pi(1+y^2)}.
\]
For $\varepsilon\in\RR$ define the secant family
\begin{equation}\label{eq:mode-secant-family}
\Psi_\varepsilon(y):=\frac{R_\varepsilon(y)-1}{\varepsilon}
=\frac{2y-\varepsilon}{1+(y-\varepsilon)^2},
\qquad y\in\RR.
\end{equation}
Since $R_\varepsilon(y)\,\omega(\dd y)=P_1(y-\varepsilon)\,\dd y$ and
$\int P_1=1$, one has $\Psi_\varepsilon\in L^2_0(\omega)$ for every
$\varepsilon\in\RR$.

\begin{theorem}[The mode Gram kernel]%
\label{thm:mode-gram-kernel}
For every $\varepsilon,\eta\in\RR$,
\begin{equation}\label{eq:mode-gram-kernel}
\langle \Psi_\varepsilon,\Psi_\eta\rangle_{L^2(\omega)}
=K(\varepsilon,\eta)
:=\frac{2}{4+(\varepsilon-\eta)^2}.
\end{equation}
\end{theorem}

\begin{proof}
Set
\[
I(\varepsilon,\eta)
:=\frac{1}{\pi}\int_{\RR}
\frac{1+y^2}{(1+(y-\varepsilon)^2)(1+(y-\eta)^2)}\,\dd y.
\]
The integrand is $O(|z|^{-2})$ on the upper half-plane with simple poles
at $z=\varepsilon+\ii$ and $z=\eta+\ii$.  A direct residue computation gives
\[
I(\varepsilon,\eta)
=\frac{\varepsilon^2+\eta^2+4}{(\varepsilon-\eta)^2+4}.
\]
Since $\int R_\varepsilon\,\dd\omega=1$, one obtains
\[
\int_{\RR}(R_\varepsilon-1)(R_\eta-1)\,\dd\omega
=I(\varepsilon,\eta)-1
=\frac{2\varepsilon\eta}{4+(\varepsilon-\eta)^2},
\]
and dividing by $\varepsilon\eta\neq 0$ gives \eqref{eq:mode-gram-kernel};
the remaining cases follow by continuity.
\end{proof}

\begin{corollary}[Closed formulas for quadratic mode integrals]%
\label{cor:mode-quadratic-integrals}
For $|\varepsilon|$ small, $\Psi_\varepsilon=\sum_{n=1}^\infty u_n\varepsilon^{n-1}$
in $L^2(\omega)$.  Comparing the power series of $K(\varepsilon,\eta)=
\frac12\sum_{\ell=0}^\infty (-1)^\ell 4^{-\ell}(\varepsilon-\eta)^{2\ell}$
with $\sum_{m,n}\langle u_m,u_n\rangle\varepsilon^{m-1}\eta^{n-1}$ gives,
for $p,q\ge 1$,
\[
\langle u_{2p},u_{2q}\rangle_{L^2(\omega)}
=(-1)^{p+q}2^{-(2p+2q-1)}\tbinom{2p+2q-2}{2p-1},
\]
for $p,q\ge 0$,
\[
\langle u_{2p+1},u_{2q+1}\rangle_{L^2(\omega)}
=(-1)^{p+q}2^{-(2p+2q+1)}\tbinom{2p+2q}{2p},
\]
and $\langle u_{2p},u_{2q+1}\rangle_{L^2(\omega)}=0$.
\end{corollary}

\begin{proposition}[The mode space and its RKHS completion]%
\label{thm:mode-space-rkhs}
Let $L^2_0(\omega):=\{f\in L^2(\omega):\int f\,\dd\omega=0\}$ and let
$\mathcal H_K$ be the RKHS with kernel $K$.  Then
$\operatorname{span}\{\Psi_\varepsilon:\varepsilon\in\RR\}$ is dense in
$L^2_0(\omega)$, and the assignment $\Psi_\varepsilon\mapsto K(\cdot,\varepsilon)$
extends to a unitary $\mathcal U:L^2_0(\omega)\to\mathcal H_K$.
\end{proposition}

\begin{proof}
Identify $L^2(\omega)$ with $L^2((-\pi/2,\pi/2),\dd\theta/\pi)$ via
$y=\tan\theta$.  Equations
\eqref{eq:cayley-mode-even-fourier}--\eqref{eq:cayley-mode-odd-fourier}
show that $\{u_n\}$ generate all trigonometric polynomials built from
$\{\cos(2k\theta),\sin(2k\theta):k\ge 1\}$, hence
$\operatorname{span}\{u_n:n\ge 1\}$ is dense in $L^2_0(\omega)$.

For $|\varepsilon|<1$ the series
$\Psi_\varepsilon=\sum_{n\ge 1}u_n\varepsilon^{n-1}$ converges in
$L^2(\omega)$, so the scalar function
$F(\varepsilon):=\langle f,\Psi_\varepsilon\rangle_{L^2(\omega)}$ is
analytic on the unit disc with Taylor coefficients $\langle f,u_n\rangle$.
If $f$ is orthogonal to all $\Psi_\varepsilon$, then $F(\varepsilon)\equiv 0$
and therefore $\langle f,u_n\rangle=0$ for every $n\ge 1$, which forces
$f=0$ by the density of $\operatorname{span}\{u_n\}$.
Hence $\operatorname{span}\{\Psi_\varepsilon\}$ is dense in $L^2_0(\omega)$.

The assignment $\Psi_\varepsilon\mapsto K(\cdot,\varepsilon)$ preserves
inner products by Theorem~\ref{thm:mode-gram-kernel}, so it extends by
density to the desired unitary $\mathcal U$.
\end{proof}

\begin{remark}[Connection to entropy asymptotics]\label{rem:rkhs-entropy}
The unitary $\mathcal U$ identifies the centred fluctuation profile
$\widetilde\delta_t(y):=h_t(\bar\gamma+ty)/g_t(\bar\gamma+ty)-1$ with a
Bochner integral in $\mathcal H_K$:
\[
(\mathcal U\widetilde\delta_t)(a)
=\frac{1}{t}\int_\RR \Delta\,K(a,\Delta/t)\,\dd\nu(\gamma),\qquad a\in\RR.
\]
This provides the functional-analytic home for $\widetilde\delta_t$ and is
consistent with the mode-expansion approach used throughout
Section~\ref{sec:precision-potential}.
\end{remark}

\begin{proposition}[RKHS as Poisson image space]%
\label{thm:gram-space-poisson-image}
The RKHS $\mathcal H_K$ is isometrically equal to
$\mathcal X:=\{f\in L^2(\RR):\int e^{2|\xi|}|\widehat f(\xi)|^2\dd\xi<\infty\}$
with norm $\|f\|_\mathcal X^2:=\frac{1}{2\pi^2}\int e^{2|\xi|}|\widehat f|^2\dd\xi$.
Equivalently, $\mathcal H_K=\sqrt\pi\,T_1(L^2(\RR))$ isometrically, and every
$f\in\mathcal H_K$ extends holomorphically to the strip
$\mathbb S:=\{z\in\CC:|\Im z|<1\}$ with pointwise bound
$|f(x+\ii y)|\le \|f\|_{\mathcal H_K}/\sqrt{2(1-y^2)}$.
\end{proposition}

\begin{proof}
The Fourier transform of $k(x):=2/(4+x^2)$ is $\widehat k(\xi)=\pi e^{-2|\xi|}$.
Reproducing: $\langle f,K(\cdot,a)\rangle_\mathcal X = f(a)$ by Fourier inversion.
Moore--Aronszajn uniqueness gives $\mathcal X=\mathcal H_K$.  The Poisson image
description and strip estimate follow by Cauchy--Schwarz as in the proof of
\cite[Proposition~3.1]{SunZhou2008}.
\end{proof}
